\subsection{Fundamento de valoración y dependencia entre riesgos}
La probabilidad ordinal describe la posibilidad de materialización durante los 56 meses del contrato, en las actividades y ventanas expuestas de cada ficha. No es una tasa mensual ni una frecuencia por maniobra. Los valores 2, 3 y 4 se fundamentan en condición limitada, mecanismo plausible sin historial suficiente y exposición recurrente o concentrada, respectivamente. La Tabla~\ref{tab:sd8-fundamentos} individualiza la causa y el criterio empleado; todos son juicios iniciales de ingeniería, sujetos a nueva evidencia. El impacto toma la consecuencia máxima de cada ficha, con escalamiento independiente de la multiplicación cuando afecta la protección obligatoria.
\begingroup\setlength{\tabcolsep}{3pt}
\begin{longtable}{L{.10\textwidth}L{.27\textwidth}L{.27\textwidth}L{.27\textwidth}}
\caption{Fundamento de las valoraciones iniciales}\label{tab:sd8-fundamentos}\\\hline
\EncabezadoTabla{ID} & \EncabezadoTabla{Condición concreta} & \EncabezadoTabla{Probabilidad propuesta} & \EncabezadoTabla{Impacto propuesto}\\\hline\endfirsthead
\hline\EncabezadoTabla{ID} & \EncabezadoTabla{Condición concreta} & \EncabezadoTabla{Probabilidad propuesta} & \EncabezadoTabla{Impacto propuesto}\\\hline\endhead
R-01 & Avisos de la autoridad marítima y ventanas costeras limitadas. & $P=4$: Hay recurrencia estacional, dependencia humana o ventana obligatoria que concentra la exposición. & $I=5$: Puede impedir un hito o afectar personas, integridad, privacidad o condición obligatoria; escalamiento por gravedad.\\\hline
R-02 & Dependencias físicas y fallos correlacionados en el recinto. & $P=3$: Existe causa y consecuencia identificables; sin historial suficiente para descartar el evento. & $I=5$: Puede impedir un hito o afectar personas, integridad, privacidad o condición obligatoria; escalamiento por gravedad.\\\hline
R-03 & Reintentos, doble escritura o reglas cloud/local incompatibles. & $P=3$: Existe causa y consecuencia identificables; sin historial suficiente para descartar el evento. & $I=5$: Puede impedir un hito o afectar personas, integridad, privacidad o condición obligatoria; escalamiento por gravedad.\\\hline
R-04 & Cliente sin TI, turnos operacionales y validación final sin holgura. & $P=4$: Hay recurrencia estacional, dependencia humana o ventana obligatoria que concentra la exposición. & $I=5$: Puede impedir un hito o afectar personas, integridad, privacidad o condición obligatoria; escalamiento por gravedad.\\\hline
R-05 & Funciones coincidentes, refuerzos no habilitados o factores de ambiente distintos del modelo. & $P=3$: Existe causa y consecuencia identificables; sin historial suficiente para descartar el evento. & $I=4$: Se comprometen habilitación, cobertura, ventana o calidad de varios resultados.\\\hline
R-06 & Fin de soporte, ficha incorrecta o componente no apto para el entorno marino. & $P=3$: Existe causa y consecuencia identificables; sin historial suficiente para descartar el evento. & $I=4$: Se comprometen habilitación, cobertura, ventana o calidad de varios resultados.\\\hline
R-07 & Servicios administrados, formatos o claves no exportables. & $P=3$: Existe causa y consecuencia identificables; sin historial suficiente para descartar el evento. & $I=4$: Se comprometen habilitación, cobertura, ventana o calidad de varios resultados.\\\hline
R-08 & Concentración operacional y frecuencia de medición mayor a la dimensionada. & $P=3$: Existe causa y consecuencia identificables; sin historial suficiente para descartar el evento. & $I=4$: Se comprometen habilitación, cobertura, ventana o calidad de varios resultados.\\\hline
R-09 & Origen no verificado, vulnerabilidad, secreto expuesto o procedencia falsificable. & $P=3$: Existe causa y consecuencia identificables; sin historial suficiente para descartar el evento. & $I=5$: Puede impedir un hito o afectar personas, integridad, privacidad o condición obligatoria; escalamiento por gravedad.\\\hline
R-10 & Permisos amplios, registros inseguros o datos reales en pruebas. & $P=3$: Existe causa y consecuencia identificables; sin historial suficiente para descartar el evento. & $I=5$: Puede impedir un hito o afectar personas, integridad, privacidad o condición obligatoria; escalamiento por gravedad.\\\hline
R-11 & Flexión, agua, combustible clasificado o interacción durante maniobra. & $P=3$: Existe causa y consecuencia identificables; sin historial suficiente para descartar el evento. & $I=5$: Puede impedir un hito o afectar personas, integridad, privacidad o condición obligatoria; escalamiento por gravedad.\\\hline
R-12 & Legado sin soporte/documentación y conocimiento concentrado. & $P=3$: Existe causa y consecuencia identificables; sin historial suficiente para descartar el evento. & $I=5$: Puede impedir un hito o afectar personas, integridad, privacidad o condición obligatoria; escalamiento por gravedad.\\\hline
R-13 & Turnos, suplencias, inglés o especialistas sin provisión efectiva. & $P=3$: Existe causa y consecuencia identificables; sin historial suficiente para descartar el evento. & $I=4$: Se comprometen habilitación, cobertura, ventana o calidad de varios resultados.\\\hline
R-14 & Personal nuevo y monitores renovados, con capacitación restringida en peak. & $P=4$: Hay recurrencia estacional, dependencia humana o ventana obligatoria que concentra la exposición. & $I=4$: Se comprometen habilitación, cobertura, ventana o calidad de varios resultados.\\\hline
R-15 & Dependencias de las innovaciones y ciclos naturales; medición tardía. & $P=3$: Existe causa y consecuencia identificables; sin historial suficiente para descartar el evento. & $I=4$: Se comprometen habilitación, cobertura, ventana o calidad de varios resultados.\\\hline
R-16 & Depósito por liberación omitido, titularidad incorrecta o custodia sin habilitar. & $P=3$: Existe causa y consecuencia identificables; sin historial suficiente para descartar el evento. & $I=4$: Se comprometen habilitación, cobertura, ventana o calidad de varios resultados.\\\hline
R-IN1-01 & Se conserva la coordinación telefónica sin ingresar sus cambios. & $P=4$: Hay recurrencia estacional, dependencia humana o ventana obligatoria que concentra la exposición. & $I=4$: Se comprometen habilitación, cobertura, ventana o calidad de varios resultados.\\\hline
R-IN1-02 & Adhesión voluntaria y pocas solicitudes elegibles. & $P=3$: Existe causa y consecuencia identificables; sin historial suficiente para descartar el evento. & $I=3$: Retrabajo o pérdida reversible de beneficio; no autoriza rebajar seguridad ni una puerta obligatoria.\\\hline
R-IN1-03 & Contratos, identidades o versiones difieren entre IN1, CTR y VAR. & $P=3$: Existe causa y consecuencia identificables; sin historial suficiente para descartar el evento. & $I=4$: Se comprometen habilitación, cobertura, ventana o calidad de varios resultados.\\\hline
R-IN1-04 & Autorización insuficiente de armador o contratista asignado. & $P=2$: Exposición condicionada a una autorización o vínculo específico; no consta fallo reiterado. & $I=5$: Puede impedir un hito o afectar personas, integridad, privacidad o condición obligatoria; escalamiento por gravedad.\\\hline
R-IN2-01 & Adhesión voluntaria o convocatoria sin grupo estable. & $P=3$: Existe causa y consecuencia identificables; sin historial suficiente para descartar el evento. & $I=4$: Se comprometen habilitación, cobertura, ventana o calidad de varios resultados.\\\hline
R-IN2-02 & Trabajo o disponibilidad no confirmado antes de la ventana. & $P=4$: Hay recurrencia estacional, dependencia humana o ventana obligatoria que concentra la exposición. & $I=4$: Se comprometen habilitación, cobertura, ventana o calidad de varios resultados.\\\hline
R-IN2-03 & Aviso marítimo o condición insegura coincidente con trabajos programados. & $P=4$: Hay recurrencia estacional, dependencia humana o ventana obligatoria que concentra la exposición. & $I=5$: Puede impedir un hito o afectar personas, integridad, privacidad o condición obligatoria; escalamiento por gravedad.\\\hline
R-IN2-04 & La coordinación preventiva depende de un servicio o contrato aún no recibido. & $P=3$: Existe causa y consecuencia identificables; sin historial suficiente para descartar el evento. & $I=4$: Se comprometen habilitación, cobertura, ventana o calidad de varios resultados.\\\hline
R-IN3-01 & Documento ilegible o formato cambiado; extractor propone valores erróneos. & $P=4$: Hay recurrencia estacional, dependencia humana o ventana obligatoria que concentra la exposición. & $I=4$: Se comprometen habilitación, cobertura, ventana o calidad de varios resultados.\\\hline
R-IN3-02 & Sesgo de automatización, formación incompleta o revisión apresurada. & $P=3$: Existe causa y consecuencia identificables; sin historial suficiente para descartar el evento. & $I=4$: Se comprometen habilitación, cobertura, ventana o calidad de varios resultados.\\\hline
R-IN3-03 & Proveedor, ubicación, retención o permiso difieren de condiciones aprobadas. & $P=3$: Existe causa y consecuencia identificables; sin historial suficiente para descartar el evento. & $I=5$: Puede impedir un hito o afectar personas, integridad, privacidad o condición obligatoria; escalamiento por gravedad.\\\hline
R-IN3-04 & Dependencia remota no disponible durante ingreso. & $P=3$: Existe causa y consecuencia identificables; sin historial suficiente para descartar el evento. & $I=3$: Retrabajo o pérdida reversible de beneficio; no autoriza rebajar seguridad ni una puerta obligatoria.\\\hline
R-IN3-05 & Volumen bajo o revisión consume el ahorro de captura. & $P=3$: Existe causa y consecuencia identificables; sin historial suficiente para descartar el evento. & $I=3$: Retrabajo o pérdida reversible de beneficio; no autoriza rebajar seguridad ni una puerta obligatoria.\\\hline
R-IN4-01 & Mezcla, dotación o mejoras ajenas cambian respecto de la base; medición no aprobada. & $P=3$: Existe causa y consecuencia identificables; sin historial suficiente para descartar el evento. & $I=5$: Puede impedir un hito o afectar personas, integridad, privacidad o condición obligatoria; escalamiento por gravedad.\\\hline
R-IN4-02 & Selección de tareas fáciles o presión para aparentar menor esfuerzo. & $P=3$: Existe causa y consecuencia identificables; sin historial suficiente para descartar el evento. & $I=5$: Puede impedir un hito o afectar personas, integridad, privacidad o condición obligatoria; escalamiento por gravedad.\\\hline
R-IN5-01 & Adhesión voluntaria, lecturas faltantes o cambios de ocupación, clima y lavado. & $P=3$: Existe causa y consecuencia identificables; sin historial suficiente para descartar el evento. & $I=4$: Se comprometen habilitación, cobertura, ventana o calidad de varios resultados.\\\hline
R-IN5-02 & Plan poco comprensible o abandono voluntario de armador/charter. & $P=3$: Existe causa y consecuencia identificables; sin historial suficiente para descartar el evento. & $I=3$: Retrabajo o pérdida reversible de beneficio; no autoriza rebajar seguridad ni una puerta obligatoria.\\\hline
R-IN5-03 & Vínculo histórico/consentimiento o autorización por embarcación no se comprueba. & $P=2$: Exposición condicionada a una autorización o vínculo específico; no consta fallo reiterado. & $I=5$: Puede impedir un hito o afectar personas, integridad, privacidad o condición obligatoria; escalamiento por gravedad.\\\hline
\end{longtable}\endgroup
\FuenteTabla{Registros generales y de innovaciones T-16; criterios de Synaptix, no frecuencias medidas.}
Los riesgos no son independientes: R-01/R-06/R-11 y R-IN2-02/03 pueden consumir la misma ventana física; R-03/R-12 comparten autoridad y conciliación; R-04/R-14 y riesgos de adopción comparten agenda funcional; R-02/R-13 comparten recuperación y cobertura. Proyecto agrupa el incidente, propaga a los consumidores y cuenta una sola suspensión o consumo. La valoración no suma productos ordinales ni multiplica probabilidades no disponibles. IN3 mantiene revisión humana aunque se estime ahorro; IN4 no mejora exposición ocultando casos abiertos o difíciles.

\subsection{Escenarios cuantificados y límites de recuperación}
Las magnitudes siguientes complementan FMEA. Se distingue presupuesto de error mensual, respuesta de incidente y RTO: cumplir cuatro horas en DR no permite gastar cuatro horas todos los meses bajo una disponibilidad crítica de 99,9\,\%. Para un mes de treinta días, $30\times24\times60\times(1-0{,}999)=43{,}2$ minutos de indisponibilidad crítica máxima; un incidente de cuatro horas excedería ese presupuesto mensual en 196,8 minutos. La medición se realiza sobre transacciones de negocio, con denominador y exclusiones contractuales, no sobre el estado del servidor. \parencite[arts.~20 y 78]{sd8-ba}

El escenario de restauración desde cero de SD5 utiliza 60 GB históricos provisionales, 6,25 MB/s de transferencia, 20 MB/s de restauración estructurada y 65 minutos de preparación, reaplicación y comprobación integrada. Se obtiene $65+60.000/(6{,}25\times60)+10.000/(20\times60)=233{,}33$ minutos; con doble volumen, 401,67 minutos. A volumen doble, conservar 6,25 MB/s no cumple cuatro horas: se requieren al menos $120.000/[60(240-65-20.000/(20\times60))]=12{,}63$ MB/s efectivos de transferencia, manteniendo los otros supuestos. Este límite calculado no acredita una red ni elimina crecimiento, tiempos de identidad o fallas de copias; se ensaya y se ajusta capacidad antes de aceptar. La conmutación desde réplica preparada utiliza otra secuencia y no se confunde con restauración desde cero.

Para mostrador, 90 legajos a 90 segundos representan 8.100 segundos, 135 minutos de trabajo serial. Cuatro ejecutores sin otras tareas reducirían el lote ideal a 33,75 minutos, pero no demostrarían la latencia individual de 90 segundos ante llegadas simultáneas ni su disponibilidad real. El análisis obliga a preparar antecedentes locales, medir cola y tiempo completo y dimensionar puestos con la demanda y distribución de llegadas; el promedio anual no prueba el peak. Si falta capacidad o autorización externa, el sistema conserva el legajo y la tarea pendiente, sin conceder zarpe. \parencite[RT-09.01; numeral 16.1, decisión 21]{sd8-c07}

La holgura de marejada es el escenario de diez días hábiles por campaña registrado en SD7, no un resultado meteorológico ni una holgura CPM demostrada. Si el margen admisible antes del hito es $M$ y la suspensión consume $s$, el déficit de calendario es $\max(0,s-M)$. Con $M=10$ días y suspensiones de 5/10/15, quedan 5/0/0 días y aparece déficit 0/0/5; son escenarios condicionados a disponer realmente de esos diez días fuera del congelamiento. Se retira del programa toda actividad suspendida, incluida remota; dos ventanas alternativas no se suman por existir en el documento. La red vigente de SD7/T-15 aún tiene duraciones y asignaciones pendientes: el plan identifica esa brecha, y no afirma cabida ni cumplimiento de hitos por restar escenarios aislados. Si el margen real es cero, cualquier suspensión positiva requiere recuperación demostrada o registra atraso, sin mover meses contractuales.

El perfil preparado de SD4 conserva 18 vCPU/72 GiB frente al peak de 32/128: el mínimo aporta sólo 56,25\,\% del perfil peak y requiere agregar 14 vCPU/56 GiB antes de admitir ese tráfico. La espera de ampliación pertenece al RTO. Mantener el mínimo es una reducción de capacidad permanente de 43,75\,\% frente a duplicar el peak, no ahorro monetario demostrado. Prueba, cuota y acceso a capacidad antes de H5/H10 controlan R-02/R-08; una promoción sin expansión y prueba de carga no se acepta.
